% GENERATED FILE. Edit canonical lessons, then run npm run generate:books.
% canonical-source-hash: fnv1a64:512a8eea44eb7c66
% canonical-lessons: RU-C337-tochka-zreniya, RU-C337-poziciya, RU-C337-vzglyad, RU-C337-zayavlyat, RU-C337-schitatsya

\chapter{My point of view}
\label{ch:ru-my-point-of-view}
\hlchaptermodality{337}

\begin{chapteropening}
\textbf{By the end of this chapter:} \emph{I can give my point of view and my position, say how I see things, state something publicly and say what is considered true.}

Complete the last lesson of chapter 337: I can give my point of view and my position, say how I see things, state something publicly and say what is considered true.
\end{chapteropening}

\section[\texorpdfstring{\ru{точка} \ru{зрения} (tóchka zréniya)}{точка зрения (tóchka zréniya)}]{\ru{точка} \ru{зрения} (tóchka zréniya) — a point of view}

\label{lesson:RU-C337-tochka-zreniya}



\begin{quote}
Before the new one: say the Russian for a task, then the Russian for experience.
\end{quote}

\subsection*{You'll want to know: \ru{точка} \ru{зрения}}
\textbf{\ru{точка} \ru{зрения}} (\emph{tóchka zréniya}) — \textquotedblleft{}a point of view\textquotedblright{}.

\begin{quote}
\ru{С} \ru{моей} \ru{точки} \ru{зрения}, \ru{это} \ru{хорошая} \ru{идея}.
\end{quote}

\begin{quote}
\emph{S moyéy tóchki zréniya, éto khoróshaya idéya.}
\end{quote}

\begin{quote}
From my point of view, it is a good idea.
\end{quote}

Feminine, after \emph{\ru{точка}}: genitive \emph{\ru{точки} \ru{зрения}}. \emph{\ru{С} \ru{моей} \ru{точки} \ru{зрения}} — from my point of view.

\begin{cousinweb}
\textbf{\ru{точка}} is usually traced to the verb \emph{\ru{ткнуть}} \textquotedblleft{}to poke, to jab\textquotedblright{}: a dot is the mark a jab leaves (EtymDB lists Lithuanian \emph{taškas} \textquotedblleft{}a dot\textquotedblright{} beside it). \textbf{\ru{зрение}} \textquotedblleft{}sight\textquotedblright{} is built on the old verb \emph{\ru{зреть}} \textquotedblleft{}to look\textquotedblright{}, so a point of view is literally the point you look from.
\end{cousinweb}

\begin{grammarlens}[title={What you've built}]
One more word for telling, explaining and describing at B1.
\end{grammarlens}

\subsection*{Your turn}
\begin{itemize}
\raggedright
  \item \emph{Say it:} \emph{tóchka zréniya}
  \item \emph{Say it:} \emph{tóchka zréniya}, once more
  \item \emph{Say it:} a sentence of your own with \emph{tóchka zréniya}
  \item \emph{Recall:} say the Russian for a task, then the Russian for experience, then say \emph{tóchka zréniya} again
\end{itemize}

\begin{tcolorbox}[breakable,colback=teal!4,colframe=teal!35!black,title={Before you move on}]
What does \ru{точка} \ru{зрения} mean? (\textquotedblleft{}A point of view\textquotedblright{}.) Say it once more.
\end{tcolorbox}
\section[\texorpdfstring{\ru{позиция} (pozítsiya)}{позиция (pozítsiya)}]{\ru{позиция} (pozítsiya) — a position, a stance}

\label{lesson:RU-C337-poziciya}



\begin{quote}
Before the new one: say the Russian for experience, then the Russian for a point of view.
\end{quote}

\subsection*{You'll want to know: \ru{позиция}}
\textbf{\ru{позиция}} (\emph{pozítsiya}) — \textquotedblleft{}a position, a stance\textquotedblright{}.

\begin{quote}
\ru{Моя} \ru{позиция} \ru{проста}: \ru{надо} \ru{больше} \ru{читать}.
\end{quote}

\begin{quote}
\emph{Moyá pozítsiya prostá: nádo bólshe chitát.}
\end{quote}

\begin{quote}
My position is simple: we need to read more.
\end{quote}

Feminine. \emph{\ru{Занимать} \ru{позицию}} — to take a position. It is also a place in a list or a team: \emph{\ru{первая} \ru{позиция}}.

\begin{cousinweb}
\textbf{\ru{позиция}} is usually traced, through western European languages, to Latin \emph{positiō} \textquotedblleft{}a placing, a position\textquotedblright{} (Lewis \& Short), which is built on \emph{pōnere} \textquotedblleft{}to put, to place\textquotedblright{}. A stance is where you have placed yourself.
\end{cousinweb}

\begin{grammarlens}[title={What you've built}]
One more word for telling, explaining and describing at B1.
\end{grammarlens}

\subsection*{Your turn}
\begin{itemize}
\raggedright
  \item \emph{Say it:} \emph{pozítsiya}
  \item \emph{Say it:} \emph{pozítsiya}, once more
  \item \emph{Say it:} a sentence of your own with \emph{pozítsiya}
  \item \emph{Recall:} say the Russian for experience, then the Russian for a point of view, then say \emph{pozítsiya} again
\end{itemize}

\begin{tcolorbox}[breakable,colback=teal!4,colframe=teal!35!black,title={Before you move on}]
What does \ru{позиция} mean? (\textquotedblleft{}A position, a stance\textquotedblright{}.) Say it once more.
\end{tcolorbox}
\section[\texorpdfstring{\ru{взгляд} (vzglyad)}{взгляд (vzglyad)}]{\ru{взгляд} (vzglyad) — a view, an opinion; a look}

\label{lesson:RU-C337-vzglyad}



\begin{quote}
Before the new one: say the Russian for a point of view, then the Russian for a position.
\end{quote}

\subsection*{You'll want to know: \ru{взгляд}}
\textbf{\ru{взгляд}} (\emph{vzglyad}) — \textquotedblleft{}a view, an opinion; a look\textquotedblright{}.

\begin{quote}
\ru{У} \ru{нас} \ru{разные} \ru{взгляды} \ru{на} \ru{жизнь}.
\end{quote}

\begin{quote}
\emph{U nas ráznyye vzglyády na zhizn.}
\end{quote}

\begin{quote}
We have different views on life.
\end{quote}

Masculine. \emph{\ru{На} \ru{мой} \ru{взгляд}} — in my view. It is also a glance: \emph{\ru{быстрый} \ru{взгляд}}.

\begin{cousinweb}
\textbf{\ru{взгляд}} is inherited from Proto-Slavic \emph{*vŭzględŭ} (EtymDB), built from the prefix \emph{*vŭz-} \textquotedblleft{}up\textquotedblright{} and \emph{*ględěti} \textquotedblleft{}to look\textquotedblright{}, the verb that survives in \emph{\ru{глядеть}}. A view began as a look upwards.
\end{cousinweb}

\begin{grammarlens}[title={What you've built}]
One more word for telling, explaining and describing at B1.
\end{grammarlens}

\subsection*{Your turn}
\begin{itemize}
\raggedright
  \item \emph{Say it:} \emph{vzglyad}
  \item \emph{Say it:} \emph{vzglyad}, once more
  \item \emph{Say it:} a sentence of your own with \emph{vzglyad}
  \item \emph{Recall:} say the Russian for a point of view, then the Russian for a position, then say \emph{vzglyad} again
\end{itemize}

\begin{tcolorbox}[breakable,colback=teal!4,colframe=teal!35!black,title={Before you move on}]
What does \ru{взгляд} mean? (\textquotedblleft{}A view, an opinion; a look\textquotedblright{}.) Say it once more.
\end{tcolorbox}
\section[\texorpdfstring{\ru{заявлять} (zayavlyát)}{заявлять (zayavlyát)}]{\ru{заявлять} (zayavlyát) — to declare, to state publicly}

\label{lesson:RU-C337-zayavlyat}



\begin{quote}
Before the new one: say the Russian for a position, then the Russian for a view.
\end{quote}

\subsection*{You'll want to know: \ru{заявлять}}
\textbf{\ru{заявлять}} (\emph{zayavlyát}) — \textquotedblleft{}to declare, to state publicly\textquotedblright{}.

\begin{quote}
\ru{Директор} \ru{заявляет}, \ru{что} \ru{завод} \ru{не} \ru{закроют}.
\end{quote}

\begin{quote}
\emph{Diréktor zayavlyáyet, chto zavód ne zakróyut.}
\end{quote}

\begin{quote}
The director states that the factory will not be closed.
\end{quote}

Imperfective; perfective \emph{\ru{заявить}} (\emph{\ru{заявлю}, \ru{заявишь}}). Formal: officials and newspapers \emph{\ru{заявляют}}, while among friends people \emph{\ru{говорят}}.

\begin{cousinweb}
\textbf{\ru{заявлять}} is built from \emph{\ru{за}-} and \emph{\ru{являть}} \textquotedblleft{}to show, to reveal\textquotedblright{}, the verb behind \emph{\ru{явный}} \textquotedblleft{}obvious\textquotedblright{}. \emph{\ru{Явить}} comes through Old East Slavic \emph{javiti} from Proto-Slavic \emph{*aviti} \textquotedblleft{}to make visible\textquotedblright{} (EtymDB): to declare is to make a thing plain.
\end{cousinweb}

\begin{grammarlens}[title={What you've built}]
One more word for telling, explaining and describing at B1.
\end{grammarlens}

\subsection*{Your turn}
\begin{itemize}
\raggedright
  \item \emph{Say it:} \emph{zayavlyát}
  \item \emph{Say it:} \emph{zayavlyát}, once more
  \item \emph{Say it:} a sentence of your own with \emph{zayavlyát}
  \item \emph{Recall:} say the Russian for a position, then the Russian for a view, then say \emph{zayavlyát} again
\end{itemize}

\begin{tcolorbox}[breakable,colback=teal!4,colframe=teal!35!black,title={Before you move on}]
What does \ru{заявлять} mean? (\textquotedblleft{}To declare, to state publicly\textquotedblright{}.) Say it once more.
\end{tcolorbox}
\section[\texorpdfstring{\ru{считаться} (schitátsya)}{считаться (schitátsya)}]{\ru{считаться} (schitátsya) — to be considered, to be regarded as}

\label{lesson:RU-C337-schitatsya}



\begin{quote}
Before the new one: say the Russian for a view, then the Russian for to declare.
\end{quote}

\subsection*{You'll want to know: \ru{считаться}}
\textbf{\ru{считаться}} (\emph{schitátsya}) — \textquotedblleft{}to be considered, to be regarded as\textquotedblright{}.

\begin{quote}
\ru{Этот} \ru{район} \ru{считается} \ru{самым} \ru{тихим} \ru{в} \ru{городе}.
\end{quote}

\begin{quote}
\emph{Étot rayón schitáyetsya sámym tíkhim v górode.}
\end{quote}

\begin{quote}
This district is considered the quietest in the city.
\end{quote}

An imperfective verb. It takes the instrumental: \emph{\ru{считается} \ru{лучшим}}. \emph{\ru{Считаться} \ru{с}} + instrumental — to take into account.

\begin{cousinweb}
\textbf{\ru{считаться}} is the reflexive of \emph{\ru{считать}}, built from \emph{\ru{с}-} and \emph{\ru{читать}}. \emph{\ru{Читать}} shares its root with \emph{\ru{число}} \textquotedblleft{}number\textquotedblright{}, from Proto-Slavic \emph{*čislo}, built on \emph{*čisti} (EtymDB): to be considered is to be counted as something.
\end{cousinweb}

\begin{grammarlens}[title={What you've built}]
One more word for telling, explaining and describing at B1.
\end{grammarlens}

\subsection*{Your turn}
\begin{itemize}
\raggedright
  \item \emph{Say it:} \emph{schitátsya}
  \item \emph{Say it:} \emph{schitátsya}, once more
  \item \emph{Say it:} a sentence of your own with \emph{schitátsya}
  \item \emph{Recall:} say the Russian for a view, then the Russian for to declare, then say \emph{schitátsya} again
\end{itemize}

\begin{tcolorbox}[breakable,colback=teal!4,colframe=teal!35!black,title={Before you move on}]
What does \ru{считаться} mean? (\textquotedblleft{}To be considered, to be regarded as\textquotedblright{}.) Say it once more.
\end{tcolorbox}
